\section{Two polynomial filters}
\label{sec:filters}

We need two different polynomials in excitation counts.  The first
approximates the zero projection on an initial part of the integer spectrum.
The second detects a tail while remaining nonnegative on the whole real
line.  Both may grow outside their effective ranges; explicit bounds on
that growth will be essential. The zero filter combines normalized
Chebyshev approximation with exact roots at the first positive integers,
a classical approximation-and-interpolation construction for symmetric
predicates; see \citet[Section~3.1 and Lemma~3.9]{Sherstov2009}.
We give the construction and its estimates explicitly because transport
requires control of the polynomial over the entire count interval,
as well as accuracy on the initial integer range.

\subsection{An elementary Chebyshev construction}

Define real polynomials $C_\ell$ by
\[
 C_0(y)=1,\qquad C_1(y)=y,\qquad
 C_{\ell+1}(y)=2yC_\ell(y)-C_{\ell-1}(y).
\]
The cosine and hyperbolic cosine addition formulas satisfy this same
recurrence and initial values, so
\[
 C_\ell(\cos\theta)=\cos(\ell\theta),
 \qquad
 C_\ell(\cosh\tau)=\cosh(\ell\tau).
\]
Thus $|C_\ell(y)|\le1$ for $-1\le y\le1$, and $C_\ell$ is positive
and nondecreasing on $[1,\infty)$.  We also have the elementary global
bound
\begin{equation}
 |C_\ell(y)|\le(2|y|+1)^\ell
 \quad(y\in\mathbb R).
 \label{eq:chebyshev-global}
\end{equation}
Indeed, writing $M=2|y|+1\ge1$, the initial cases hold, and the
recurrence gives
$2|y|M^\ell+M^{\ell-1}\le M^{\ell+1}$ because
$2|y|M+1\le M^2$.

For $1\le a<R\le m$, define
\[
 y(x)=1+\frac{2(a-x)}{R-a},\qquad
 S_\ell(x)=\frac{C_\ell(y(x))}{C_\ell(y(0))}.
\]
The denominator is a positive real constant.  As $x$ increases from zero
to $a$, $y(x)$ decreases from $y(0)>1$ to one; as $x$ increases from $a$
to $R$, $y(x)$ decreases from one to minus one.  The preceding formulas
therefore imply
\begin{equation}
 \begin{aligned}
 S_\ell(0)&=1,\\
 |S_\ell(x)|&\le1 &&(0\le x\le R),\\
 |S_\ell(x)|&\le C_\ell(y(0))^{-1} &&(a\le x\le R).
 \end{aligned}
 \label{eq:normalized-chebyshev}
\end{equation}
Whenever $a\le R/2$, we have $R-a\ge1$.  For every $x\in[0,m]$ it
follows that $|y(x)|\le1+2m\le3m$.  Since $C_\ell(y(0))\ge1$,
\eqref{eq:chebyshev-global} yields the uniform bound
\begin{equation}
 \sup_{0\le x\le m}|S_\ell(x)|\le(7m)^\ell.
 \label{eq:normalized-global}
\end{equation}

For a useful exact expression for the denominator, put $z=\sqrt{a/R}$.
Defining
\[
 \operatorname{artanh}z=\frac12\log\frac{1+z}{1-z}
 \quad(0\le z<1),
\]
direct substitution into the exponential formula for $\cosh$ gives
\[
 y(0)=\frac{1+z^2}{1-z^2}
      =\cosh\bigl(2\operatorname{artanh}z\bigr).
\]
Moreover, integration of $(1-w^2)^{-1}$ between zero and $z$ gives
\begin{equation}
 z\le\operatorname{artanh}z\le\frac{z}{1-z^2}.
 \label{eq:artanh-bound}
\end{equation}
Consequently, if $\tau_0=2\operatorname{artanh}\sqrt{a/R}$, then
\[
 C_\ell(y(0))=\cosh(\ell\tau_0),\qquad
 \tau_0\ge2\sqrt{a/R},\qquad
 \tau_0=(2+o(1))\sqrt{a/R}
\]
as $a/R\to0$.  These facts prove all the approximation estimates below.

\subsection{Approximating the zero projection}

\begin{lemma}[Integer zero filter]
\label{lem:zero-filter}
Fix $0<q<L<1$ and $r>(q+L)/2$.  For all sufficiently large integers $m$
there is a real polynomial $Q$ such that
\[
 \deg Q\le m^r,\qquad
 \sup_{0\le x\le m}|Q(x)|\le\exp(m^r),
\]
and, for every integer $j$ with $0\le j\le m^L$,
\[
 \bigl|Q(j)-\one_{j=0}\bigr|\le\exp(-m^q).
\]
In fact the construction has $Q(0)=1$ exactly.
\end{lemma}

\begin{proof}
Set
\[
 a=\lceil m^q\rceil,\qquad R=m^L,\qquad
 \ell=\left\lceil10\sqrt{aR}\log(2m)\right\rceil.
\]
Eventually $a\le R/2$, since $q<L$.  The denominator estimate above and
$\cosh x\ge e^x/2$ for $x\ge0$ give
\[
 C_\ell(y(0))
 \ge\frac12\exp\bigl(2\ell\sqrt{a/R}\bigr)
 \ge\frac12(2m)^{20a}
 \ge(2m)^{5a}.
\]
Define
\[
 Q(x)=S_\ell(x)\prod_{i=1}^{a}\left(1-\frac{x}{i}\right).
\]
Then $Q(0)=1$, and $Q$ vanishes at each integer from one through $a$.
For $0\le x\le m$, the product has absolute value at most $(m+1)^a$.
For $a\le x\le R$, \eqref{eq:normalized-chebyshev} therefore gives
\[
 |Q(x)|\le\frac{(m+1)^a}{(2m)^{5a}}
       \le(2m)^{-4a}\le\exp(-m^q).
\]
This proves the required accuracy at all the integer arguments.  Accuracy
between the integers in $(0,a)$ is neither claimed nor needed.

For the degree and global size, put $\beta=(q+L)/2$.  Since $q<\beta$,
\[
 \deg Q\le\ell+a=O\bigl(m^\beta\log(2m)\bigr).
\]
The estimate \eqref{eq:normalized-global} holds on the entire real
interval $[0,m]$, including between the small integer roots, and gives
\[
 \sup_{0\le x\le m}|Q(x)|\le(7m)^\ell(m+1)^a.
\]
In particular,
\[
 \log\sup_{0\le x\le m}|Q(x)|
 \le\ell\log(7m)+a\log(m+1)
 =O\bigl(m^\beta(\log(2m))^2\bigr).
\]
For any fixed $r>\beta$, both the degree and this logarithmic bound are
at most $m^r$ eventually.  Indeed, every fixed power of $\log(2m)$ is
$o(m^{r-\beta})$, as is seen by setting $w=\log m$ and comparing a
polynomial in $w$ with $e^{(r-\beta)w}$.
\end{proof}

Since every $H_A$ has integer spectrum, this is a projection
approximation in the required sense.  More explicitly, with
$\Pi=\one_{H_A\le m^L}$, the spectral theorem gives
\[
 \norm{\bigl(P_A-Q(H_A)\bigr)\Pi}\le e^{-m^q},
 \qquad \norm{Q(H_A)}\le e^{m^r}.
\]
For every vector $\xi$, splitting into the two spectral ranges yields
\begin{equation}
 \norm{\bigl(P_A-Q(H_A)\bigr)\xi}
 \le e^{-m^q}\norm{\xi}
    +(1+e^{m^r})\norm{\one_{H_A>m^L}\xi}.
 \label{eq:zero-filter-vector}
\end{equation}
Its operator degree is at most $m^r$ by
Lemma~\ref{lem:degree-shift}.

\subsection{A nonnegative tail detector}

\begin{lemma}[Tail filter]
\label{lem:tail-filter}
Fix $0<v<h<u\le1$ and $s>0$.  Let $k$ be a degree exponent satisfying
$k>s+(u-h)/2$.  For all sufficiently large integers $m$, there is a real
polynomial $T$ with
\[
 \deg T\le m^k,\qquad
 T(x)\ge0\quad(x\in\mathbb R),\qquad
 \sup_{0\le x\le m}T(x)\le\exp(m^k),
\]
and
\begin{align*}
 0\le T(x)&\le1 &&(0\le x\le m^u),\\
 T(x)&\le\exp(-m^s) &&(0\le x\le m^v),\\
 T(x)&\ge\tfrac14 &&(m^h\le x\le m^u).
\end{align*}
\end{lemma}

\begin{proof}
Set
\[
 a=m^h,\qquad R=m^u,\qquad
 \ell=\left\lceil3\sqrt{R/a}\right\rceil,
 \qquad B(x)=1-S_\ell(x)^2.
\]
Eventually $a\le R/2$, and \eqref{eq:normalized-chebyshev} gives
$0\le B\le1$ on $[0,R]$.  The denominator satisfies
$C_\ell(y(0))\ge\cosh6$ by \eqref{eq:artanh-bound}.  Hence, on $[a,R]$,
\[
 B(x)\ge1-(\cosh6)^{-2}>\frac78.
\]
For example, $\cosh6\ge1+6^2/2=19$ proves the last strict inequality.

We next verify uniformly that $B(x)\to0$ on $[0,m^v]$.
Put $z=\sqrt{a/R}$, which tends to zero, and for this low range put
\[
 z_x=\sqrt{\frac{a-x}{R-x}}.
\]
These are positive eventually because $m^v/a\to0$.
The same hyperbolic identity gives
\[
 y(x)=\cosh\bigl(2\operatorname{artanh}z_x\bigr),
 \qquad
 S_\ell(x)=
 \frac{\cosh(2\ell\operatorname{artanh}z_x)}
      {\cosh(2\ell\operatorname{artanh}z)}.
\]
Now
\[
 \frac{z_x}{z}
 =\sqrt{\frac{1-x/a}{1-x/R}}\longrightarrow1
 \quad\text{uniformly for }0\le x\le m^v,
 \qquad \ell z\longrightarrow3.
\]
Also $0<z_x\le z$ in this range.  The inequalities
\eqref{eq:artanh-bound} show uniformly that
$\operatorname{artanh}z_x/z_x\to1$ and
$\operatorname{artanh}z/z\to1$.  Both hyperbolic cosine arguments
therefore tend uniformly to six.  Their ratio tends uniformly to one,
which proves $0\le B(x)\le1/8$ on $[0,m^v]$ eventually.

Choose an odd integer $J$ with $10m^s\le J<10m^s+2$ and define the
majority polynomial
\[
 M_J(z)=\sum_{i=(J+1)/2}^{J}\binom Ji z^i(1-z)^{J-i}.
\]
For $z\in[0,1]$ its summands are nonnegative and their full binomial
sum is one, so $0\le M_J(z)\le1$.  On $0\le z\le1/8$,
\[
 M_J(z)\le2^J z^{J/2}\le2^{-J/2}.
\]
Changing the summation index in the complementary binomial sum gives
$M_J(1-z)=1-M_J(z)$, since $J$ is odd.  Therefore
$M_J(z)\ge1-2^{-J/2}$ for $7/8\le z\le1$.

Take
\[
 T(x)=\bigl(M_J(B(x))\bigr)^2.
\]
This square of a real polynomial is nonnegative at every real argument.
It lies in $[0,1]$ on $[0,R]$.  On $[0,m^v]$ it is at most
$2^{-J}\le\exp(-m^s)$, since $J\ge10m^s$ and $10\log2>1$.
On $[a,R]$ it is at least $(1-2^{-J/2})^2\ge1/4$ eventually.

It remains to control its degree and its growth on $[0,m]$, where
$B$ need not lie in $[0,1]$.  Write $A=(7m)^{2\ell}\ge1$.
By \eqref{eq:normalized-global}, throughout this whole interval,
\[
 |B(x)|\le1+A,\qquad |1-B(x)|\le A.
\]
Taking absolute values in the binomial formula, and then including all
indices in its sum, gives
\[
 |M_J(B(x))|
 \le\bigl(|B(x)|+|1-B(x)|\bigr)^J
 \le(1+2A)^J\le(3A)^J.
\]
Thus the required global estimates, including the final square, are
\[
 \deg T\le4\ell J,\qquad
 \sup_{0\le x\le m}|T(x)|\le9^J(7m)^{4\ell J}.
\]
With $\beta=s+(u-h)/2$, we have $\ell J=O(m^\beta)$, and hence
\[
 \deg T=O(m^\beta),\qquad
 \log\sup_{0\le x\le m}|T(x)|
 \le J\log9+4\ell J\log(7m)
 =O\bigl(m^\beta\log(2m)\bigr).
\]
The fixed strict inequality $k>\beta$ makes each bound at most $m^k$
eventually.  This proves the lemma, including the endpoint case $u=1$.
\end{proof}

Evaluating $T$ at a count operator gives a positive semidefinite operator
of degree at most $m^k$ and norm at most $e^{m^k}$.  In particular, its
global nonnegativity prevents contributions from outside $[m^h,m^u]$
from cancelling the detected tail in a quadratic form.
